%%=============================================================================
%% Inleiding
%%=============================================================================

\chapter{\IfLanguageName{dutch}{Inleiding}{Introduction}}%
\label{ch:inleiding}

%De inleiding moet de lezer net genoeg informatie verschaffen om het onderwerp te begrijpen en in te zien waarom de onderzoeksvraag de moeite waard is om te onderzoeken. In de inleiding ga je literatuurverwijzingen beperken, zodat de tekst vlot leesbaar blijft. Je kan de inleiding verder onderverdelen in secties als dit de tekst verduidelijkt. Zaken die aan bod kunnen komen in de inleiding~\autocite{Pollefliet2011}:

%\begin{itemize}
%  \item context, achtergrond
%  \item afbakenen van het onderwerp
%  \item verantwoording van het onderwerp, methodologie
%  \item probleemstelling
%  \item onderzoeksdoelstelling
%  \item onderzoeksvraag
%  \item \ldots
%\end{itemize}

De afgelopen jaren is de digitalisering binnen lokale besturen sterk toegenomen. Gemeenten maken steeds meer gebruik van digitale systemen om administratieve processen efficiënter te laten verlopen. Tegelijkertijd neemt de werkdruk binnen overheidsorganisaties toe, onder andere doordat het takenpakket van lokale besturen blijft uitbreiden. Dit zorgt ervoor dat administratieve processen steeds complexer worden en vaak verschillende diensten met elkaar moeten samenwerken. \newline

Een belangrijk administratief proces binnen organisaties is het onboardingproces van nieuwe medewerkers. Tijdens dit proces moeten verschillende stappen worden uitgevoerd om ervoor te zorgen dat een nieuwe medewerker toegang krijgt tot de nodige systemen, applicaties en infrastructuur. Binnen een gemeentebestuur betekent dit bijvoorbeeld het aanmaken van gebruikersaccounts, het instellen van toegangsrechten, het activeren van badge-toegang tot gebouwen en het registreren van medewerkers in verschillende interne systemen. Omdat meerdere diensten hierbij betrokken zijn, is een goede coördinatie essentieel. \newline


\section{\IfLanguageName{dutch}{Probleemstelling}{Problem Statement}}%
\label{sec:probleemstelling}

%Uit je probleemstelling moet duidelijk zijn dat je onderzoek een meerwaarde heeft voor een concrete doelgroep. De doelgroep moet goed gedefinieerd en afgelijnd zijn. Doelgroepen als ``bedrijven,'' ``KMO's'', systeembeheerders, enz.~zijn nog te vaag. Als je een lijstje kan maken van de personen/organisaties die een meerwaarde zullen vinden in deze bachelorproef (dit is eigenlijk je steekproefkader), dan is dat een indicatie dat de doelgroep goed gedefinieerd is. Dit kan een enkel bedrijf zijn of zelfs één persoon (je co-promotor/opdrachtgever).

Het huidige onboardingproces binnen het gemeentebestuur van Berlare verloopt grotendeels manueel en vereist tussenkomst van meerdere diensten en systemen.Hierdoor ontstaat een proces dat tijdrovend en foutgevoelig is. Informatie moet vaak manueel worden ingevoerd of doorgegeven tussen verschillende medewerkers, wat het risico op miscommunicatie of vertragingen verhoogt. Ook is het huidige proces op basis van een Excelbestand dat niet meer up to date is. Hierdoor gaat informatie verloren of worden er overbodige gegevens meegegeven.\newline

De primaire doelgroep van dit onderzoek zijn de medewerkers van het gemeentebestuur van Berlare die betrokken zijn bij het onboardingproces van nieuwe werknemers. Dit omvat onder andere medewerkers van de personeelsdienst en de IT-dienst. Voor deze groep kan een efficiënter onboardingproces leiden tot minder administratieve lasten, minder fouten en een betere samenwerking tussen diensten.\newline

Indirect kan ook de dienstverlening aan inwoners van de gemeente baat hebben bij een efficiëntere interne werking. Wanneer interne processen vlotter verlopen, kunnen medewerkers sneller en efficiënter hun taken uitvoeren, wat uiteindelijk de kwaliteit van de gemeentelijke dienstverlening kan verbeteren.

\section{\IfLanguageName{dutch}{Onderzoeksvraag}{Research question}}%
\label{sec:onderzoeksvraag}

%Wees zo concreet mogelijk bij het formuleren van je onderzoeksvraag. Een onderzoeksvraag is trouwens iets waar nog niemand op dit moment een antwoord heeft (voor zover je kan nagaan). Het opzoeken van bestaande informatie (bv. ``welke tools bestaan er voor deze toepassing?'') is dus geen onderzoeksvraag. Je kan de onderzoeksvraag verder specifiëren in deelvragen. Bv.~als je onderzoek gaat over performantiemetingen, dan 

Op basis van de vastgestelde problematiek wordt in deze bachelorproef onderzocht hoe het onboardingproces binnen het gemeentebestuur kan worden geoptimaliseerd door middel van automatisatie. \newline

De centrale onderzoeksvraag luidt als volgt:

Op welke manier kan een geautomatiseerd systeem ontwikkeld worden dat, op basis van gegevens uit de personeelsdienst, een volledige en correcte onboarding kan uitvoeren binnen de IT-infrastructuur van het gemeentebestuur van Berlare? \newline

Om deze onderzoeksvraag te beantwoorden worden een aantal deelvragen onderzocht:
\begin{itemize}
    \item Hoe verloopt het huidige onboardingproces binnen het gemeentebestuur van Berlare?
    \item Welke stappen en systemen zijn momenteel betrokken bij dit proces?
    \item Welke technologieën en systemen bestaan er die automatisatie van onboardingprocessen mogelijk maken?
    \item Op welke manier kan een geautomatiseerde oplossing geïntegreerd worden binnen de bestaande IT-infrastructuur?
    \item Hoeveel tijdswinst is er ten opzichte van het huidige onboardingsproces
    \item Aan welke vereisten moet de oplossing voldoen
\end{itemize}
    

\section{\IfLanguageName{dutch}{Onderzoeksdoelstelling}{Research objective}}%
\label{sec:onderzoeksdoelstelling}

%Wat is het beoogde resultaat van je bachelorproef? Wat zijn de criteria voor succes? Beschrijf die zo concreet mogelijk. Gaat het bv.\ om een proof-of-concept, een prototype, een verslag met aanbevelingen, een vergelijkende studie, enz.

Het doel van deze bachelorproef is het ontwikkelen van een proof-of-concept of prototype dat het onboardingproces gedeeltelijk of volledig kan automatiseren. In dit prototype moet een medewerker van de personeelsdienst het onboardingproces kunnen starten, waarna de verdere stappen automatisch worden uitgevoerd binnen de verschillende betrokken systemen. \newline

Het succes van het onderzoek wordt beoordeeld aan de hand van verschillende criteria. Zo moet het prototype aantonen dat het onboardingproces efficiënter kan verlopen dan in de huidige situatie. Daarnaast moet de oplossing het aantal manuele handelingen verminderen en het risico op fouten verkleinen. Tot slot zal een vergelijking worden gemaakt tussen de huidige tijdsbesteding van het onboardingproces en de tijd die nodig is bij gebruik van de voorgestelde oplossing.

\section{\IfLanguageName{dutch}{Opzet van deze bachelorproef}{Structure of this bachelor thesis}}%
\label{sec:opzet-bachelorproef}

% Het is gebruikelijk aan het einde van de inleiding een overzicht te
% geven van de opbouw van de rest van de tekst. Deze sectie bevat al een aanzet
% die je kan aanvullen/aanpassen in functie van je eigen tekst.

De rest van deze bachelorproef is als volgt opgebouwd:

In Hoofdstuk~\ref{ch:stand-van-zaken} wordt een overzicht gegeven van de stand van zaken binnen het onderzoeksdomein, op basis van een literatuurstudie.

In Hoofdstuk~\ref{ch:methodologie} wordt de methodologie toegelicht en worden de gebruikte onderzoekstechnieken besproken om een antwoord te kunnen formuleren op de onderzoeksvragen.

% TODO: Vul hier aan voor je eigen hoofstukken, één of twee zinnen per hoofdstuk

In Hoofdstuk~\ref{ch:Huidig-Onboardingsproces} wordt het huidige onboardingsproces besproken. Hierbij komen de valkuilen aan bod die er momenteel zijn. 

In Hoofdstuk~\ref{ch:Scripting} worden de powershell scripts besproken die gebruikt worden om het onboardingsproces te automatiseren. 

In Hoofdstuk~\ref{ch:Power Automate} worden de flows in Power Automate besproken. Deze flows vormen samen met de scripts uit ~\ref{ch:Scripting} het automatische onboardingsproces.

In Hoofdstuk~\ref{ch:conclusie}, tenslotte, wordt de conclusie gegeven en een antwoord geformuleerd op de onderzoeksvragen. Daarbij wordt ook een aanzet gegeven voor toekomstig onderzoek binnen dit domein.